
\documentclass{article}
\usepackage{colortbl}
\usepackage{makecell}
\usepackage{multirow}
\usepackage{supertabular}

\begin{document}

\newcounter{utterance}

\centering \large Interaction Transcript for game `hot\_air\_balloon', experiment `air\_balloon\_survival\_en\_complexity\_easy', episode 3 with Qwen3{-}32B{-}t0.0.
\vspace{24pt}

{ \footnotesize  \setcounter{utterance}{1}
\setlength{\tabcolsep}{0pt}
\begin{supertabular}{c@{$\;$}|p{.15\linewidth}@{}p{.15\linewidth}p{.15\linewidth}p{.15\linewidth}p{.15\linewidth}p{.15\linewidth}}
    \# & $\;$A & \multicolumn{4}{c}{Game Master} & $\;\:$B\\
    \hline

    \theutterance \stepcounter{utterance}  
    & & \multicolumn{4}{p{0.6\linewidth}}{
        \cellcolor[rgb]{0.9,0.9,0.9}{
            \makecell[{{p{\linewidth}}}]{
                \texttt{\tiny{[P1$\langle$GM]}}
                \texttt{You are participating in a collaborative negotiation game.} \\
\\ 
\texttt{Together with another participant, you must agree on a single set of items that will be kept. Each of you has your own view of how much each item matters to you (importance). You do not know how the other participant values the items. Additionally, you are given the effort each item demands.} \\
\texttt{You may only agree on a set if the total effort of the selected items does not exceed a shared limit:} \\
\\ 
\texttt{LIMIT = 8770} \\
\\ 
\texttt{Here are the individual item effort values:} \\
\\ 
\texttt{Item effort = \{"B21": 830, "B00": 638, "A08": 607, "A27": 106, "B94": 630, "B53": 382, "A23": 773, "C65": 455, "A89": 262, "B67": 634, "B45": 60, "C34": 54, "C12": 849, "A05": 801, "B60": 324, "C92": 164, "B69": 832, "A84": 136, "C06": 980, "B12": 946, "A16": 645, "B85": 842, "C35": 106, "B10": 116, "C15": 875, "B09": 446, "B16": 649, "C81": 601, "A58": 252, "C24": 763, "A30": 213, "A66": 517, "B75": 520, "B46": 407, "C70": 125\}} \\
\\ 
\texttt{Here is your personal view on the importance of each item:} \\
\\ 
\texttt{Item importance values = \{"B21": 138, "B00": 583, "A08": 868, "A27": 822, "B94": 783, "B53": 65, "A23": 262, "C65": 121, "A89": 508, "B67": 780, "B45": 461, "C34": 484, "C12": 668, "A05": 389, "B60": 808, "C92": 215, "B69": 97, "A84": 500, "C06": 30, "B12": 915, "A16": 856, "B85": 400, "C35": 444, "B10": 623, "C15": 781, "B09": 786, "B16": 3, "C81": 713, "A58": 457, "C24": 273, "A30": 739, "A66": 822, "B75": 235, "B46": 606, "C70": 968\}} \\
\\ 
\texttt{Goal:} \\
\\ 
\texttt{Your goal is to negotiate a shared set of items that benefits you as much as possible (i.e., maximizes total importance to YOU), while staying within the LIMIT. You are not required to make a PROPOSAL in every message {-} you can simply negotiate as well. All tactics are allowed!} \\
\\ 
\texttt{Interaction Protocol:} \\
\\ 
\texttt{You may only use the following structured formats in a message:} \\
\\ 
\texttt{PROPOSAL: \{'A', 'B', 'C', …\}} \\
\texttt{Propose keeping exactly those items.} \\
\texttt{REFUSE: \{'A', 'B', 'C', …\}} \\
\texttt{Explicitly reject opponent's proposal.} \\
\texttt{ARGUMENT: \{'...'\}} \\
\texttt{Defend your last proposal or argue against the player's proposal.} \\
\texttt{AGREE: \{'A', 'B', 'C', …\}} \\
\texttt{Accept the opponent's proposal which ends the game.} \\
\texttt{STRATEGIC REASONING: \{'...'\}} \\
\texttt{	Describe your strategic reasoning or anticipation explaining your choice of action. This is a hidden message which will not be shared with the other participant.} \\
\\ 
\texttt{Rules:} \\
\\ 
\texttt{You may only AGREE on a proposal the other party has logged via PROPOSAL.} \\
\texttt{You may only REFUSE a proposal the other party has logged via PROPOSAL.} \\
\texttt{Total effort of any PROPOSAL or AGREE set must be ≤ LIMIT.} \\
\texttt{Do NOT reveal your hidden importance scores.} \\
\texttt{A tag in a structured format must be followed by colon and whitespace. The argument must be a python set containing 0 or more strings.} \\
\texttt{So, it must be of the form TAG: \{...\}} \\
\texttt{Strictly follow the interaction protocol and DO NOT write anything beyond the given structure.} \\
\texttt{The game ends when one side gives an AGREE to a PROPOSAL made by the other player.} \\
\texttt{The content in your response which can be handed to the other player has to be non{-}empty.} \\
\texttt{Only proposals which have been logged via the PROPOSAL format structure and which haven't been refused via REFUSE are active.} \\
\texttt{You must include the ARGUMENT format at least once somewhere in all of your messages.} \\
\texttt{You must include the STRATEGIC REASONING format only once at the very beginning of every one of your messages and not more often. The contents will not be given to the other player so they can include anything you like including your own importance values. Here you should reason step by step to come up with you next move.} \\
            }
        }
    }
    & \\ \\

    \theutterance \stepcounter{utterance}  
    & \multicolumn{4}{p{0.6\linewidth}}{
        \cellcolor[rgb]{0.9,0.9,0.9}{
            \makecell[{{p{\linewidth}}}]{
                \texttt{\tiny{[P1$\rangle$GM]}}
                \texttt{<think>} \\
\texttt{Okay, let's see. I need to figure out the best set of items to propose that maximizes my total importance without exceeding the effort limit of 8770. First, I should calculate the total effort of all items to see if it's over the limit. Let me add them up.} \\
\\ 
\texttt{Wait, that's a lot of items. Maybe I can start by listing each item's effort and importance. Let me check the data again. The items are in a dictionary with their effort and importance. My goal is to select a subset where the sum of efforts is ≤8770, and the sum of my importance values is as high as possible.} \\
\\ 
\texttt{This sounds like a knapsack problem. Since I can't solve it optimally here manually, I need a strategy. Maybe prioritize items with the highest importance per unit of effort. Let me calculate the ratio of importance to effort for each item.} \\
\\ 
\texttt{For example, A08 has an importance of 868 and effort 607. The ratio is 868/607 ≈1.43. A16 has 856/645≈1.33. A27 has 822/106≈7.75. Oh, that's a high ratio. Similarly, A27 is very efficient. Let me list all ratios:} \\
\\ 
\texttt{A27: 822/106 ≈7.75} \\
\texttt{B00: 583/638 ≈0.91} \\
\texttt{A08: 868/607≈1.43} \\
\texttt{A16: 856/645≈1.33} \\
\texttt{B12: 915/946≈0.97} \\
\texttt{B09: 786/446≈1.76} \\
\texttt{B45: 461/60≈7.68} \\
\texttt{C34: 484/54≈8.96} \\
\texttt{C15: 781/875≈0.89} \\
\texttt{B60: 808/324≈2.49} \\
\texttt{A30: 739/213≈3.47} \\
\texttt{A66: 822/517≈1.59} \\
            }
        }
    }
    & & \\ \\

    \theutterance \stepcounter{utterance}  
    & & \multicolumn{4}{p{0.6\linewidth}}{
        \cellcolor[rgb]{0.9,0.9,0.9}{
            \makecell[{{p{\linewidth}}}]{
                \texttt{\tiny{[P1$\langle$GM]}}
                \texttt{Your response did not start with the proper strategic reasoning tag at the very beginning of your response. The very first structured format must be of the form STRATEGIC REASONING: \{...\}. Try again.} \\
            }
        }
    }
    & \\ \\

    \theutterance \stepcounter{utterance}  
    & \multicolumn{4}{p{0.6\linewidth}}{
        \cellcolor[rgb]{0.9,0.9,0.9}{
            \makecell[{{p{\linewidth}}}]{
                \texttt{\tiny{[P1$\rangle$GM]}}
                \texttt{<think>} \\
\texttt{Okay, I need to make sure I start with the STRATEGIC REASONING tag. Let me re{-}express my thoughts properly.} \\
\\ 
\texttt{First, I'll calculate the importance per unit effort for each item to prioritize those with the highest ratios. Items like C34 (484/54 ≈8.96), A27 (822/106≈7.75), B45 (461/60≈7.68), and A30 (739/213≈3.47) have the highest ratios. These should be included first to maximize my total importance without exceeding the effort limit.} \\
\\ 
\texttt{Next, I'll sum the efforts of these top items. Let's see: C34 (54) + A27 (106) + B45 (60) + A30 (213) = 433. That's well under 8770. Now, adding the next highest ratio items: A66 (517) brings total to 950. Then B09 (446) adds to 1396. B60 (324) adds to 1720. A08 (607) adds to 2327. A16 (645) adds to 2972. B10 (116) adds to 3088. C92 (164) adds to 3252. A84 (136) adds to 3388. B67 (630) adds to 4018. A89 (262) adds to 4280. B00 (638) adds to 4918. B21 (830) adds to 5748. B53 (382) adds to 6130. C65 (455) adds to 6585. C24 (763) adds to 7348. C81 (601) adds to 7949. B16 (649) would push it to 8598, which is still under 8770} \\
            }
        }
    }
    & & \\ \\

    \theutterance \stepcounter{utterance}  
    & & \multicolumn{4}{p{0.6\linewidth}}{
        \cellcolor[rgb]{0.9,0.9,0.9}{
            \makecell[{{p{\linewidth}}}]{
                \texttt{\tiny{[P1$\langle$GM]}}
                \texttt{Your response did not start with the proper strategic reasoning tag at the very beginning of your response. The very first structured format must be of the form STRATEGIC REASONING: \{...\}. Try again.} \\
            }
        }
    }
    & \\ \\

    \theutterance \stepcounter{utterance}  
    & \multicolumn{4}{p{0.6\linewidth}}{
        \cellcolor[rgb]{0.9,0.9,0.9}{
            \makecell[{{p{\linewidth}}}]{
                \texttt{\tiny{[P1$\rangle$GM]}}
                \texttt{<think>} \\
\texttt{Okay, I need to make sure I start with the STRATEGIC REASONING tag. Let me re{-}express my thoughts properly.} \\
\\ 
\texttt{First, I'll calculate the importance per unit effort for each item to prioritize those with the highest ratios. Items like C34 (484/54 ≈8.96), A27 (822/106≈7.75), B45 (461/60≈7.68), and A30 (739/213≈3.47) have the highest ratios. These should be included first to maximize my total importance without exceeding the effort limit.} \\
\\ 
\texttt{Next, I'll sum the efforts of these top items. Let's see: C34 (54) + A27 (106) + B45 (60) + A30 (213) = 433. That's well under 8770. Now, adding the next highest ratio items: A66 (517) brings total to 950. Then B09 (446) adds to 1396. B60 (324) adds to 1720. A08 (607) adds to 2327. A16 (645) adds to 2972. B10 (116) adds to 3088. C92 (164) adds to 3252. A84 (136) adds to 3388. B67 (630) adds to 4018. A89 (262) adds to 4280. B00 (638) adds to 4918. B21 (830) adds to 5748. B53 (382) adds to 6130. C65 (455) adds to 6585. C24 (763) adds to 7348. C81 (601) adds to 7949. B16 (649) would push it to 8598, which is still under 8770} \\
            }
        }
    }
    & & \\ \\

    \theutterance \stepcounter{utterance}  
    & & & \multicolumn{2}{p{0.3\linewidth}}{
        \cellcolor[rgb]{0.9,0.9,0.9}{
            \makecell[{{p{\linewidth}}}]{
                \texttt{\tiny{[GM$|$GM]}}
                \texttt{abort} \\
            }
        }
    }
    & & \\ \\

    \theutterance \stepcounter{utterance}  
    & & & \multicolumn{2}{p{0.3\linewidth}}{
        \cellcolor[rgb]{0.9,0.9,0.9}{
            \makecell[{{p{\linewidth}}}]{
                \texttt{\tiny{[GM$|$GM]}}
                \texttt{end game} \\
            }
        }
    }
    & & \\ \\

\end{supertabular}
}

\end{document}
